\chapter*{Preface}
\markboth{Preface}{Preface}
Welcome to \textit{Azure Master Data Engineer}. This text is designed as an academic, hands-on, progressive program for building cloud-scale data platforms on Microsoft Azure. By working through the architecture discussions, command-line labs, and design exercises, you will move from storage foundations to governed lakehouse, warehouse, streaming, and machine-learning data products.

This book is the theory-and-architecture companion to a six-milestone Azure Data Engineering curriculum. Where the practice workspace asks you to build weekly deliverables, this book explains why each Azure service exists, what operational failure mode it addresses, and how a data engineer reasons about reliability, cost, compliance, and maintainability.

\section*{Who This Book Is For}
This book is written for data engineers, analytics engineers, cloud engineers, database professionals, and technically minded analysts preparing for Azure data-platform work. It assumes basic SQL, Python, and command-line familiarity. Prior Azure experience is helpful but not required: the early chapters introduce storage accounts, identity, naming, redundancy, and lifecycle controls before the program moves into Synapse, Data Factory, Databricks, Event Hubs, Cosmos DB, Azure Machine Learning, and governance.

\section*{How This Book Is Organized}
The book follows the training roadmap as \textbf{6 Parts / 6 Milestones} and \textbf{24 Chapters / 24 Weeks}, followed by appendices:
\begin{itemize}
    \item \textbf{Part I: Core Data Infrastructure and Storage} establishes ADLS Gen2, secure access, file formats, medallion zones, and governance foundations.
    \item \textbf{Part II: Data Warehousing and Analytics} builds analytical serving layers with Azure Synapse Analytics, serverless SQL, dedicated SQL pools, and Power BI semantic models.
    \item \textbf{Part III: Batch Processing and ETL Orchestration} develops repeatable pipelines with Azure Data Factory, integration runtimes, Databricks, Spark, and Delta Lake.
    \item \textbf{Part IV: Streaming and Real-Time Analytics} introduces Event Hubs, IoT ingestion, Stream Analytics, real-time Spark, and operational analytics patterns with Cosmos DB.
    \item \textbf{Part V: Machine Learning and MLOps} connects curated data products to Azure Machine Learning, feature engineering, pipeline automation, model monitoring, and GenAI-ready data services.
    \item \textbf{Part VI: Security, Governance, and Exam Prep} integrates Zero Trust, observability, FinOps, compliance, DP-203-style exam review, and the enterprise capstone.
\end{itemize}
Each chapter opens with \textbf{Learning Objectives} and closes with \textbf{Further Reading}, \textbf{Key Takeaways}, and \textbf{Exercises}. Hands-on chapters include Azure CLI, JSON, SQL, or Python listings that are meant to be adapted in a real subscription.

\section*{Reading Paths}
\begin{itemize}
    \item \textbf{New to Azure data engineering:} read Part I in order, perform every hands-on storage exercise, then continue to Synapse and Data Factory.
    \item \textbf{Experienced SQL or BI practitioner:} skim Chapter 1 for storage vocabulary, then read Parts II and III while returning to governance notes as needed.
    \item \textbf{Spark or Databricks practitioner:} read Chapters 1--4 for Azure storage and governance assumptions, then focus on Databricks, Delta Lake, streaming, and MLOps.
    \item \textbf{Platform architect:} prioritize architecture diagrams, pitfall boxes, Well-Architected trade-offs, and the security/governance chapters.
    \item \textbf{Certification candidate:} complete the conceptual exercises, review the command-line examples, and use Part VI to consolidate DP-203 objectives.
\end{itemize}

\section*{Conventions}
Code identifiers, resource names, and file paths appear in \texttt{monospace}. Azure CLI snippets use the shared listing style and assume PowerShell on Windows unless noted. Callout boxes flag \textbf{objectives}, \textbf{key terms}, \textbf{definitions}, \textbf{notes}, \textbf{tips}, and \textbf{pitfalls}. Citations refer to the bibliography at the back of the book. Service names use their current Microsoft branding, while design guidance emphasizes transferable principles: least privilege, explicit recovery objectives, lifecycle economics, reproducible infrastructure, observability, and governed data contracts.
